Considera la funció $f(x)=x^3-4x$.

\begin{apartats}

\apartat[1,25]{0,75}
Troba els punts on la gràfica de $f$ talla l'eix $X$, i una primitiva de $f$.

\begin{solucio}
$x^3-4x=x(x-2)(x+2)=0$: talla l'eix a $x=-2$, $x=0$ i $x=2$. Una primitiva és
$F(x)=\dfrac{x^4}4-2x^2$.
\end{solucio}

\begin{tria}{area-signe}
\itemtria{area}{1}{1,25}
Calcula l'àrea de la regió tancada per la gràfica de $f$ i l'eix $X$.

\begin{solucio}
Hi ha dues regions, $[-2,0]$ i $[0,2]$, i cal integrar-les per separat:
\[
\left|\int_{-2}^0 f(x)\,dx\right|+\left|\int_0^2 f(x)\,dx\right|
=\bigl|F(0)-F(-2)\bigr|+\bigl|F(2)-F(0)\bigr|=|0-(4-8)|+|(4-8)-0|=4+4=8.
\]
L'àrea és de 8 unitats quadrades.
\end{solucio}

\itemtria{error-signe}{1}{1,25}
Un company calcula $\displaystyle\int_{-2}^2 f(x)\,dx$, obté 0 i diu que l'àrea tancada per la gràfica i
l'eix $X$ és 0. Explica on és l'error i calcula l'àrea correcta.

\begin{solucio}
La integral és correcta: $F(2)-F(-2)=-4-(-4)=0$. Però a $[-2,0]$ la funció és positiva, i a $[0,2]$,
negativa: la integral hi dona $+4$ i $-4$, que es compensen. L'àrea és la suma dels valors absoluts de
les integrals de cada tros: $|4|+|-4|=8$ unitats quadrades.
\end{solucio}
\end{tria}

\begin{nomesllarg}
\apartat{0,75}
Calcula l'àrea de la regió tancada per la paràbola $y=-x^2+4$ i l'eix $X$.

\begin{solucio}
Talla l'eix a $x=\pm2$, i entre aquests valors és positiva:
$\displaystyle\int_{-2}^2(-x^2+4)\,dx=\left[-\frac{x^3}3+4x\right]_{-2}^2=\frac{16}3-\left(-\frac{16}3\right)=\frac{32}3$
unitats quadrades.
\end{solucio}
\end{nomesllarg}

\end{apartats}
